
% !TEX TS-program = pdflatex
\documentclass[11pt]{article}

\usepackage[margin=1in]{geometry}
\usepackage{amsmath,amssymb,amsthm}
\usepackage{mathtools}
\usepackage{microtype}
\usepackage{enumitem}
\usepackage{hyperref}

\hypersetup{
  colorlinks=true,
  linkcolor=blue,
  citecolor=blue,
  urlcolor=blue
}

\title{Grounded Cognitive Infrastructure for Agents and ATP:\\
World-Model Calculus, Probabilistic Logic, Behavioral Contracts, and Transfer\\
(AIrxiv Primer, 2026-02-11)}
\author{AIrxiv note (conceptual synthesis)}
\date{2026-02-11}

\newcommand{\Prob}{\mathbb{P}}
\newcommand{\E}{\mathbb{E}}
\newcommand{\KL}{\mathrm{KL}}
\newcommand{\Odds}{\mathrm{Odds}}

\begin{document}
\maketitle

\begin{abstract}
This note distills a set of practical, research-oriented ideas for building neuro-symbolic systems
without skipping the basics.
The focus is on \emph{grounded incrementality}: start from strong public baselines, build evaluation
harnesses early, and introduce novel mechanisms only when they measurably improve outcomes.

We develop a coherent ``stack'' connecting:
(1) tool-using LLM agents as fast implementation and exploration engines,
(2) a \emph{world-model calculus} view of probabilistic reasoning (carry posterior state; query views),
(3) Probabilistic Logic Networks (PLN) grounded via classical baselines (Naive Bayes, k-NN, Bayesian/factor graphs),
(4) inference control as decision-making under uncertainty (bandits, forward/backward scheduling),
(5) behavioral contracts derived from operational semantics (OSLF) as a clean proposal/certification split,
and (6) transfer mechanisms (``TransWeave'') treated as typed maps with certificates and negative-transfer detection.

The intent is strictly conceptual: a primer for future design discussions and for setting up reproducible experiments.
\paragraph{Status note.} Implementation-status and count claims in this draft are snapshot-scoped; verify against current repository trackers and build outputs before citing as current state. Lean-proven claims are only those tied to explicit Lean file/theorem references; broader mathematical or historical claims are literature-backed context and not asserted as Lean-proven unless explicitly stated.
\end{abstract}

\tableofcontents

\section{Grounded incrementality: the meta-strategy}

Ambitious cognitive-architecture work often fails in one of two ways:
\begin{itemize}[leftmargin=2em]
\item \textbf{Vision without traction:} many conceptual modules, few measurable wins.
\item \textbf{Traction without coherence:} strong bottom-up systems (e.g., coding agents) that work,
but whose improvements are hard to accumulate, audit, or steer over long horizons.
\end{itemize}

A grounded synthesis strategy is:

\begin{enumerate}[label=\textbf{G\arabic*.}, leftmargin=2em]
\item \textbf{Use strong baselines as ``semantic anchors.''} If a proposed reasoning layer cannot
reproduce well-understood baselines on controlled tasks, its semantics are probably drifting.
\item \textbf{Separate semantics from control.} Semantics: what is represented and what is true under the model.
Control: what gets attention, which inference step to try next, how budgets are allocated.
\item \textbf{Make assumptions explicit.} Fast procedures are correct only under declared model classes (e.g., conditional independence).
\item \textbf{Treat ``novel'' mechanisms as optional power-ups.} They should plug in behind stable interfaces,
and must earn their complexity by measurable gains.
\end{enumerate}

This immediately suggests a concrete tactic: pick a domain where baselines are mature and evaluation is objective.
Automated theorem proving (ATP) with premise selection is such a domain.

\section{LLM agents as ``hands'': accelerate the build, keep the audit trail}

Tool-using LLM agents (especially coding agents) are powerful because they compress implementation cycles.
However, they do not automatically provide:
\begin{itemize}[leftmargin=2em]
\item stable long-term memory with provenance,
\item reliable multi-step reasoning under strict correctness constraints,
\item safety boundaries for tool access,
\item reproducible evaluation across iterations.
\end{itemize}

A productive division of labor is:
\begin{itemize}[leftmargin=2em]
\item \textbf{LLMs:} generate candidates (code, hypotheses, plans, rule sketches), do rapid iteration, propose fixes.
\item \textbf{Structured substrate:} store state explicitly, track provenance, enforce invariants, and support verifiable reasoning kernels.
\end{itemize}

This is not ``LLM as mere natural-language interface.'' It is \emph{LLM as a high-capacity proposal engine},
paired with a system that \emph{records, checks, and composes} what is proposed.

\section{World-model calculus: carry posterior state; query views}

A repeated confusion in uncertain reasoning is mixing three layers:
\begin{enumerate}[leftmargin=2em]
\item \textbf{World-model state} (what is carried and updated): a posterior representation of uncertainty,
including correlations, provenance, and context.
\item \textbf{Evidence calculus} (how new observations update state): combination rules for support and refutation.
\item \textbf{Query views} (what is reported): marginals, conditionals, scores, rankings.
\end{enumerate}

Slogan:
\begin{quote}
\emph{Carry the posterior; query the marginals.}
\end{quote}

This matters because many ``truth-value'' approaches store only local summaries (e.g., a scalar per link) and
expect global reasoning to remain valid. In general, correlations make this impossible without additional structure.
A world-model calculus discipline keeps the semantics honest: the state is the posterior; reported ``truth values''
are derived views.

\section{Evidence pairs and Naive Bayes: a minimal grounding bridge}

\subsection{Evidence algebra}

For binary hypotheses, represent evidence as an \emph{unnormalized} pair
\[
e = \langle e^+, e^- \rangle,\quad e^+,e^- \ge 0,
\]
interpretable as masses for $H$ and $\neg H$.

Define:
\[
\mathrm{strength}(e) = \frac{e^+}{e^+ + e^-},
\qquad
\mathrm{odds}(e) = \frac{e^+}{e^-}.
\]

Let evidence composition be coordinatewise multiplication:
\[
\langle a^+,a^- \rangle \otimes \langle b^+,b^- \rangle
= \langle a^+b^+, a^-b^- \rangle.
\]

\subsection{Naive Bayes reduction}

Naive Bayes for a binary hypothesis $H$ with features $x_1,\dots,x_n$ gives:
\[
\Prob(H \mid x_1,\dots,x_n)
=
\frac{\Prob(H)\prod_i \Prob(x_i\mid H)}
{\Prob(H)\prod_i \Prob(x_i\mid H) + \Prob(\neg H)\prod_i \Prob(x_i\mid \neg H)}.
\]

Encode:
\[
e_{\mathrm{prior}}=\langle \Prob(H),\Prob(\neg H)\rangle,\quad
e_i=\langle \Prob(x_i\mid H),\Prob(x_i\mid \neg H)\rangle.
\]
Then
\[
e_{\mathrm{post}} = e_{\mathrm{prior}} \otimes \bigotimes_i e_i
=
\left\langle
\Prob(H)\prod_i \Prob(x_i\mid H),\;
\Prob(\neg H)\prod_i \Prob(x_i\mid \neg H)
\right\rangle,
\]
and
\[
\Prob(H\mid x_1,\dots,x_n) = \mathrm{strength}(e_{\mathrm{post}}).
\]

This is a strategic anchor: any ``PLN-like'' evidence mechanism should be able to reproduce this exact behavior
under the Naive Bayes assumptions.

\subsection{A crucial empirical lesson: negative examples}

In premise selection (and many relevance tasks), training on positives only is typically weak.
A pos+neg formulation (binary classification) is vastly stronger because it learns what \emph{does not} predict usefulness.
This is one of the simplest ``grounded wins'' to bake into an experimental ladder.

\section{Premise selection for ATP: a clean experimental sandbox}

\subsection{Task framing}

Let $U(P,C)\in\{0,1\}$ mean ``premise $P$ is useful for proving conjecture $C$''.
Given features $F(C)$ extracted from $C$, premise selection estimates:
\[
\Prob\big(U(P,C)=1\mid F(C)\big),
\]
then ranks premises by this score and passes the top-$N$ to a prover.

\subsection{k-NN and NB as complementary baselines}

\textbf{k-NN}:
retrieve nearest past conjectures in feature space and aggregate their used premises.
k-NN often behaves like retrieval/activation spreading; it is best viewed as \emph{proposal/control}.

\textbf{Naive Bayes}:
treat premise usefulness as a Bernoulli label with feature likelihoods.
NB is fast, surprisingly strong, and gives a clean probabilistic semantics under explicit assumptions.

A practical integration pattern:
\begin{itemize}[leftmargin=2em]
\item k-NN produces a shortlist of candidate premises (cheap retrieval).
\item NB (or a richer model) reranks within the shortlist (semantic scoring).
\end{itemize}

\subsection{Repository hygiene and reproducibility}

ATP experiments accumulate artifacts quickly: splits, features, models, proofs, logs.
A robust pattern is:
\begin{itemize}[leftmargin=2em]
\item \textbf{Splits} stored as versioned JSON with fixed random seeds.
\item \textbf{Experiment folders} keyed by date/ID containing config, outputs, and logs.
\item \textbf{A single results table} summarizing all runs (metrics + parameters).
\end{itemize}

This is not bureaucracy; it is how experiments compound rather than becoming a pile of irreproducible demos.

\subsection{Data plumbing as a hidden risk}

A practical lesson from ATP work is that dataset formats and prover toolchains can dominate effort:
FOF vs THF, bushy vs chainy, external axiom includes, typed numerics, and prover flags can make or break a pipeline.
The conceptual point is: \emph{infrastructure choices are part of the scientific method} when evaluation depends on them.
Treat them as first-class, versioned components.

\section{Beyond NB: factor graphs and message passing}

Naive Bayes assumes conditional independence of features given usefulness.
A principled extension is to introduce a factor graph (or Bayesian network) that can represent dependencies:
\begin{itemize}[leftmargin=2em]
\item \textbf{Feature factors} linking $U(P,C)$ to observed features.
\item \textbf{Premise--premise factors} capturing co-use: if $P$ is used, $Q$ becomes more likely.
\end{itemize}

Inference then becomes message passing (belief propagation / variational approximations).
This yields a controlled ladder:
\begin{enumerate}[leftmargin=2em]
\item Match NB exactly under NB assumptions.
\item Add explicit dependency factors and measure net gains.
\item Keep the approximation mode explicit (tree exactness vs loopy heuristics).
\end{enumerate}

\subsection{Relation to GNNs and activation spreading}

Graph neural networks and activation-spreading heuristics can be viewed as different points in the same design space:
\begin{itemize}[leftmargin=2em]
\item Both propagate signals over graphs.
\item GNNs learn propagation operators end-to-end.
\item Factor graphs specify propagation from a probabilistic semantics (when the model class is declared).
\end{itemize}

A grounded strategy is to treat learned propagation as a proposal/ranking method whose outputs are wrapped as
evidence factors (next section), allowing hybrid systems that remain auditable.

\section{Inference control as decision-making under uncertainty}

The main practical barrier in symbolic/probabilistic reasoning is combinatorial explosion.
Inference control is the art of deciding \emph{what to try next} under budgets.

A simple, grounded framing:
\begin{itemize}[leftmargin=2em]
\item Each candidate action (apply a rule, expand a premise, run a test) has a cost.
\item Each action has an uncertain expected value (progress, information gain).
\end{itemize}

This suggests bandit-style control (e.g., Thompson sampling, UCB) for exploration/exploitation.
A more geometric framing (sometimes called ``geodesic control'') treats control as balancing:
\begin{itemize}[leftmargin=2em]
\item \textbf{Forward feasibility:} what the current state makes reachable.
\item \textbf{Backward desirability:} what moves toward goal satisfaction.
\end{itemize}

When formalized as path-space regularization, this often resembles KL-regularized control or Schr\"odinger-bridge-like objectives.
The key discipline remains: start with simple bandit/utility schedulers and upgrade only if the richer objective provides
measurable improvements.

\section{Integrating ML into reasoning without category errors}

A productive (and testable) notion of ``ML as reasoning'' is:

\begin{quote}
\emph{ML modules produce semantically typed evidence or structure proposals that are consumed by a world-model calculus.}
\end{quote}

Three canonical wrappers:
\begin{enumerate}[leftmargin=2em]
\item \textbf{Likelihood wrapper:} a model output becomes an evidence factor for a query variable.
\item \textbf{Structure wrapper:} a learner proposes a graph/rule structure class $\Sigma$ (e.g., factor edges),
which is then recorded as an assumption for inference.
\item \textbf{Proposal wrapper:} a learner proposes candidate objects (premises/rules/plans) that are later certified.
\end{enumerate}

This avoids the trap of pretending that training is deduction.
Instead, training contributes to the posterior state, and reasoning queries that state.

\subsection{Introspection at scale: probes, not global scans}

Exposing internal products of large models (embeddings, attention summaries) to a symbolic substrate does \emph{not}
mean ``inspect everything.'' It means:
\begin{itemize}[leftmargin=2em]
\item define a small set of probes/heads/readouts,
\item store only top-$K$ relations or compressed summaries,
\item treat introspection as budgeted and contextual.
\end{itemize}

This keeps the interface scalable while enabling structured memory and auditing.

\section{OSLF and native types: behavioral contracts from operational semantics}

Operational Semantics in Logical Form (OSLF) derives modal operators from a rewrite relation:
\[
\Diamond \varphi\ \text{(``may reach $\varphi$'')},\qquad
\Box \psi\ \text{(``all predecessors satisfy $\psi$'')},
\]
with a Galois connection $\Diamond \dashv \Box$.

A key practical interpretation:
\begin{itemize}[leftmargin=2em]
\item \textbf{Proposal side ($\Diamond$):} generate possible actions/plans/steps.
\item \textbf{Certification side ($\Box$):} enforce invariants, rely conditions, safety boundaries.
\end{itemize}

Because these ``types'' are behavioral properties, general checking can be verification-hard.
The practical path is decidable fragments and bounded checking: sound-but-incomplete certifications that still deliver value.

This lens is especially useful for agent systems:
LLMs propose, the certification layer gates tool access and memory updates.

\section{Transfer (``TransWeave''): keep it thin, typed, and testable}

Transfer across tasks and representations is vital for compounding capability.
However, ``transfer'' becomes hand-wavy unless represented as an object with interfaces.

A grounded formulation treats transfer as:
\begin{itemize}[leftmargin=2em]
\item \textbf{Transfer maps} between artifact types (features, rules, options, embeddings),
\item \textbf{Certificates} describing invariants preserved or tests passed,
\item \textbf{Negative-transfer flags} (explicit refusal conditions),
\item \textbf{Order-effect measurement} for composing transfers.
\end{itemize}

This makes transfer engineering-friendly: it is a module you can test, not a metaphysical leap.
Only after the thin layer works should one introduce more ambitious geometry/optimal-transport interpretations.

\section{Representation: embeddings, graphs, and ``contradictions''}

A recurring tension:
\begin{itemize}[leftmargin=2em]
\item Embeddings are compact but blur structure.
\item Explicit structure can be combinatorial (tensor-product blowups, large graphs).
\end{itemize}

A pragmatic compromise:
\begin{enumerate}[leftmargin=2em]
\item Use embeddings for similarity and retrieval.
\item Use graph/hypergraph structure for explicit compositional constraints.
\item Use context and provenance to manage apparent contradictions.
\end{enumerate}

Many ``contradictions'' are actually:
\begin{itemize}[leftmargin=2em]
\item \textbf{Context mismatch} (different implicit conditions),
\item \textbf{Compression aliasing} (distinct structures mapped to the same representation),
\item or \textbf{genuine inconsistency} requiring explicit management (paraconsistent bookkeeping).
\end{itemize}

Extensions such as complex-amplitude ``interference'' can be treated as optional diagnostics for multi-path blending,
but should not be on the critical path until the probabilistic core is grounded.

\section{Deliver value early: build a client-facing wedge}

A practical hedge for any long-horizon AGI program is to deliver value early via a narrowly-scoped wedge.
Two strong candidates that preserve the AGI vision:
\begin{itemize}[leftmargin=2em]
\item \textbf{ATP premise selection} with reproducible baselines and clear metrics.
\item \textbf{Agent verification loops:} hypothesis/evidence ledgers that reduce repeated failures and make decisions auditable.
\end{itemize}

A minimal agent-facing world-model service can expose:
assert (claim + evidence + provenance), query (belief + confidence), contradictions, explain (top evidence), recommend-next-actions (under budget).
Such a module provides immediate practical utility while serving as a scaffold for richer world-model calculus and inference control.

\section{A concrete ladder of progress}

A plausible incremental ladder that stays grounded:

\begin{enumerate}[label=\textbf{R\arabic*.}, leftmargin=2em]
\item Reproduce baseline ATP results (no learning).
\item Add k-NN and NB premise selection (pos+neg) and log net gains/losses.
\item Implement NB scoring via evidence algebra; match NB exactly (semantics anchor).
\item Add k-NN shortlist + evidence rerank; measure.
\item Introduce factor-graph reranking with co-use edges; measure.
\item Add bandit-style inference control; measure compute vs performance.
\item Introduce behavioral certification gates (OSLF-style) for safe tool use; measure reduced regressions.
\item Add transfer maps with certificates; measure faster adaptation across domains.
\end{enumerate}

Each rung is independently useful, testable, and publishable.

\section{Informal pointers (non-exhaustive)}

This note is conceptual; the following are useful entry points:
\begin{itemize}[leftmargin=2em]
\item Premise selection and ATP: k-NN/NB baselines and later learning-based selectors (e.g., MaSh; Mizar/MPTP premise selection literature).
\item Probabilistic graphical models: Bayesian networks, factor graphs, belief propagation.
\item Behavioral specifications and modal logics: OSLF and ``native types'' approaches (types as behavioral predicates; Galois connections).
\item Compositional distributional semantics: tensor-based composition and categorical formalisms for meaning.
\item Bandit/RL control and ``grain of truth'' intuition in multi-agent learning: model-class adequacy for robust learning.
\end{itemize}

\end{document}
